\documentclass[11pt,a4paper]{article}
\usepackage{get-uebung}

\blatt{5}{Spannungsteiler und Stromteiler}{Kap.~2.3--2.4}

\begin{document}

\blattkopf{
  \item Teilspannungen in Reihenschaltungen direkt mit der Spannungsteilerregel berechnen,
  \item einen Spannungsteiler für eine gewünschte Ausgangsspannung dimensionieren,
  \item erklären, warum eine Last den Spannungsteiler verfälscht, und den Fehler abschätzen,
  \item ein belastetes Potentiometer berechnen,
  \item Teilströme mit der Stromteilerregel -- mit Widerständen und mit Leitwerten -- bestimmen.
}

% -------------------------------------------------------------------
\begin{aufgabe}{Unbelasteter Spannungsteiler}{12}
\gegeben{$U = \SI{12}{V}$, $R_1 = \SI{1}{k\ohm}$, $R_2 = \SI{2}{k\ohm}$, $R_3 = \SI{3}{k\ohm}$}
\begin{center}
\begin{circuitikz}
  \draw (0,0) to[V, v<=$U$] (0,4.5) -- (3,4.5)
        to[R, l_=$R_1$, v^=$U_1$] (3,3)
        to[R, l_=$R_2$, v^=$U_2$] (3,1.5)
        to[R, l_=$R_3$, v^=$U_3$] (3,0) -- (0,0);
\end{circuitikz}
\end{center}
\begin{teilaufgaben}
  \item Berechnen Sie $U_1$, $U_2$ und $U_3$ mit der Spannungsteilerregel, \emph{ohne} vorher den Strom zu berechnen.
  \item Wie groß ist die Spannung über $R_2$ und $R_3$ zusammen? Lösen Sie dies mit einer einzigen Spannungsteiler-Gleichung.
  \item Entwurf: Aus $U = \SI{12}{V}$ sollen mit zwei Widerständen $R_a$ (oben) und $R_b$ (unten) genau $\SI{5}{V}$ an $R_b$ erzeugt werden. Durch den Teiler soll ein Querstrom von $\SI{1}{mA}$ fließen. Bestimmen Sie $R_a$ und $R_b$.
\end{teilaufgaben}
\end{aufgabe}

\begin{loesung}
\schritt{a) Spannungsteilerregel}
\leitfrage{Warum darf hier die Spannungsteilerregel verwendet werden?}
Durch alle drei Widerstände fließt derselbe Strom (keine Abzweigung). Dann verhalten sich die Spannungen wie die Widerstände, mit $R_\text{ges} = \SI{6}{k\ohm}$:
\[
  U_1 = U\frac{R_1}{R_\text{ges}} = \SI{12}{V}\cdot\frac{1}{6} = \ergebnis{\SI{2}{V}},\quad
  U_2 = \SI{12}{V}\cdot\frac{2}{6} = \ergebnis{\SI{4}{V}},\quad
  U_3 = \SI{12}{V}\cdot\frac{3}{6} = \ergebnis{\SI{6}{V}}
\]
Kontrolle mit dem Maschensatz: $2 + 4 + 6 = 12$ \checkmark
\schritt{b) Mehrere Widerstände zusammenfassen}
Im Zähler steht einfach die Summe der Widerstände, an denen die gesuchte Spannung abfällt:
\[ U_{23} = U\frac{R_2+R_3}{R_1+R_2+R_3} = \SI{12}{V}\cdot\frac{5}{6} = \ergebnis{\SI{10}{V}} \]
\schritt{c) Dimensionierung}
\leitfrage{Welche Gleichung liefert der Querstrom, welche das Teilerverhältnis?}
Aus dem Querstrom folgt der Gesamtwiderstand: $R_a + R_b = U/I = \SI{12}{V}/\SI{1}{mA} = \SI{12}{k\ohm}$.\\
Aus der Spannungsteilerregel: $\frac{R_b}{R_a+R_b} = \frac{\SI{5}{V}}{\SI{12}{V}} \;\Rightarrow\; R_b = \frac{5}{12}\cdot\SI{12}{k\ohm} = \ergebnis{\SI{5}{k\ohm}}$, \quad $R_a = \ergebnis{\SI{7}{k\ohm}}$.
\typfehler{Im Zähler den „falschen“ Widerstand einsetzen. Beim \emph{Spannungs}teiler steht im Zähler der Widerstand, an dem die gesuchte Spannung abfällt -- beim Stromteiler (Aufgabe 5.4) ist es genau umgekehrt!}
\end{loesung}

% -------------------------------------------------------------------
\begin{aufgabe}{Belasteter Spannungsteiler}{18}
Ein Spannungsteiler aus $R_1 = R_2 = \SI{10}{k\ohm}$ liegt an $U = \SI{12}{V}$. An $R_2$ wird ein Verbraucher $R_L$ angeschlossen.
\begin{center}
\begin{circuitikz}
  \draw (0,0) to[V, v<=$U$] (0,3) -- (2.5,3)
        to[R, l_=$R_1$, v^=$U_1$] (2.5,1.5)
        to[R, l_=$R_2$, -*] (2.5,0) -- (0,0);
  \draw (2.5,1.5) node[circ]{} -- (5,1.5) to[R, l_=$R_L$] (5,0) -- (2.5,0);
  \draw (5,1.5) to[short, *-o] (6.5,1.5);
  \draw (5,0) to[short, *-o] (6.5,0);
  \draw (6.5,1.5) to[open, v^=$U_2$] (6.5,0);
\end{circuitikz}
\end{center}
\begin{teilaufgaben}
  \item Wie groß ist $U_2$ ohne Last ($R_L \to \infty$)?
  \item Berechnen Sie $U_2$ für $R_L = \SI{10}{k\ohm}$, $\SI{100}{k\ohm}$ und $\SI{1}{M\ohm}$ sowie jeweils die relative Abweichung vom Leerlaufwert.
  \item Ein Kommilitone rechnet bei Last weiterhin mit $U_2 = U\cdot\frac{R_2}{R_1+R_2}$. Erklären Sie mit Hilfe der Gültigkeitsbedingung der Spannungsteilerregel, warum das falsch ist.
  \item Leiten Sie eine Faustregel ab: Wie groß muss $R_L$ im Vergleich zu $R_2$ sein, damit der Fehler unter etwa $\SI{5}{\percent}$ bleibt?
\end{teilaufgaben}
\end{aufgabe}

\begin{loesung}
\schritt{a) Leerlauf}
$U_2 = U\frac{R_2}{R_1+R_2} = \SI{12}{V}\cdot\frac{10}{20} = \ergebnis{\SI{6}{V}}$
\schritt{b) Mit Last -- erst zusammenfassen, dann teilen}
\leitfrage{Durch welche Elemente fließt jetzt derselbe Strom?}
Durch $R_1$ fließt der Gesamtstrom, der sich danach auf $R_2$ und $R_L$ aufteilt. Fassen wir $R_2\parallel R_L$ zusammen, entsteht wieder ein unverzweigter Teiler aus $R_1$ und $R_2\parallel R_L$:
\[ U_2 = U\cdot\frac{R_2\parallel R_L}{R_1 + R_2\parallel R_L} \]
\begin{center}
\begin{tabular}{cccc}
\toprule
$R_L$ & $R_2\parallel R_L$ & $U_2$ & Abweichung vom Leerlauf\\\midrule
$\SI{10}{k\ohm}$  & $\SI{5}{k\ohm}$     & $\SI{4.00}{V}$ & $\SI{-33.3}{\percent}$\\
$\SI{100}{k\ohm}$ & $\SI{9.09}{k\ohm}$  & $\SI{5.71}{V}$ & $\SI{-4.8}{\percent}$\\
$\SI{1}{M\ohm}$   & $\SI{9.90}{k\ohm}$  & $\SI{5.97}{V}$ & $\SI{-0.5}{\percent}$\\\bottomrule
\end{tabular}
\end{center}
Beispielrechnung für die erste Zeile: $U_2 = \SI{12}{V}\cdot\frac{5}{10+5} = \SI{4}{V}$.
\schritt{c) Warum die einfache Formel versagt}
Die Spannungsteilerregel setzt voraus, dass durch $R_1$ und $R_2$ \emph{derselbe} Strom fließt. Mit Last fließt durch $R_1$ der Strom $I_2 + I_L$, durch $R_2$ nur $I_2$. Die Bedingung ist verletzt. Erst nach Zusammenfassen von $R_2\parallel R_L$ gilt sie wieder.
\schritt{d) Faustregel}
Aus der Tabelle: Bei $R_L = 10\cdot R_2$ beträgt der Fehler knapp $\SI{5}{\percent}$, bei $R_L = 100\cdot R_2$ etwa $\SI{0.5}{\percent}$.
\[ \ergebnis{R_L \gtrsim 10\cdot R_2} \]
Merke: Je kleiner der Querstrom durch den Teiler im Vergleich zum Laststrom, desto stärker bricht die Spannung ein.
\typfehler{Den Lastwiderstand in Reihe statt parallel zu $R_2$ einrechnen. Und: Die Spannung kann durch eine Last nur \emph{sinken}, nie steigen -- ein größeres Ergebnis als im Leerlauf ist immer ein Rechenfehler.}
\end{loesung}

% -------------------------------------------------------------------
\begin{aufgabe}{Belastetes Potentiometer}{12}
Ein Potentiometer mit $R = \SI{10}{k\ohm}$ liegt an $U_0 = \SI{10}{V}$. Der Schleifer teilt es in die Teile $(1-p)R$ und $pR$ mit $0 \le p \le 1$. Zwischen Schleifer und unterem Ende ist eine Last $R_L = \SI{10}{k\ohm}$ angeschlossen.
\begin{center}
\begin{circuitikz}
  \draw (-0.6,0) to[V, v<=$U_0$] (-0.6,3.5) -- (2.5,3.5)
        to[R, l_=$(1-p)R$] (2.5,1.75)
        to[R, l_=$pR$] (2.5,0) -- (-0.6,0);
  \draw[dashed, gray] (0.7,-0.2) rectangle (3.3,3.7);
  \node[gray, font=\small] at (2,3.95) {Potentiometer};
  \draw (2.5,1.75) to[short, *-] (5,1.75) to[R, l_=$R_L$, v^=$U$] (5,0) -- (2.5,0) node[circ]{};
\end{circuitikz}
\end{center}
\begin{teilaufgaben}
  \item Geben Sie $U$ in Abhängigkeit von $p$ für den unbelasteten Fall ($R_L\to\infty$) an.
  \item Berechnen Sie $U$ mit Last für die Mittelstellung $p = 0{,}5$.
  \item Bei welcher Schleiferstellung $p$ liegen mit Last genau $\SI{5}{V}$ an $R_L$?
\end{teilaufgaben}
\end{aufgabe}

\begin{loesung}
\schritt{a) Unbelastet}
Beide Teile werden vom selben Strom durchflossen, also gilt die Spannungsteilerregel:
\[ U = U_0\frac{pR}{(1-p)R + pR} = \ergebnis{p\cdot U_0} \]
Die Ausgangsspannung ist linear in $p$ -- genau dafür benutzt man ein Potentiometer.
\schritt{b) Mit Last bei $p = 0{,}5$}
Unterer Teil parallel zur Last: $pR\parallel R_L = \SI{5}{k\ohm}\parallel\SI{10}{k\ohm} = \SI{3.33}{k\ohm}$.
\[ U = U_0\frac{pR\parallel R_L}{(1-p)R + pR\parallel R_L} = \SI{10}{V}\cdot\frac{3{,}33}{5 + 3{,}33} = \ergebnis{\SI{4}{V}} \]
Statt $\SI{5}{V}$ wie im Leerlauf -- die Last „zieht die Spannung herunter“.
\schritt{c) Gesuchte Schleiferstellung}
\leitfrage{$U = U_0/2$ bedeutet, dass beide Teile des Teilers gleich groß sein müssen. Welche zwei Widerstände sind das?}
Bedingung: $(1-p)R = pR\parallel R_L$. Mit $R = R_L$ (alle Werte in $\si{k\ohm}$):
\[
  10(1-p) = \frac{10p\cdot 10}{10p + 10} = \frac{10p}{p+1}
  \;\Rightarrow\; (1-p)(1+p) = p
  \;\Rightarrow\; p^2 + p - 1 = 0
\]
\[ p = -\tfrac12 + \sqrt{\tfrac14 + 1} = \frac{\sqrt5 - 1}{2} \approx \ergebnis{0{,}618} \]
(die negative Lösung ist physikalisch sinnlos). Probe: $pR\parallel R_L = 6{,}18\parallel 10 \approx \SI{3.82}{k\ohm} = (1-p)R$ \checkmark
\typfehler{Mit Last weiter $U = pU_0$ rechnen. Bei c) die Gleichung ohne Hilfsidee aufstellen -- das geht auch, wird aber länglich. Die Überlegung „gleiche Teilspannungen $\Leftrightarrow$ gleiche Widerstände“ spart viel Rechnung.}
\end{loesung}

% -------------------------------------------------------------------
\begin{aufgabe}{Stromteiler mit Widerständen und Leitwerten}{15}
Ein Strom $I = \SI{60}{mA}$ teilt sich auf drei parallele Widerstände auf: $R_1 = \SI{20}{\ohm}$, $R_2 = \SI{30}{\ohm}$, $R_3 = \SI{60}{\ohm}$.
\begin{center}
\begin{circuitikz}
  \draw (0,0) to[I, i=$I$] (0,3) -- (2,3) -- (6,3);
  \draw (2,3) to[R, l_=$R_1$, i>^=$I_1$, *-*] (2,0);
  \draw (4,3) to[R, l_=$R_2$, i>^=$I_2$, *-*] (4,0);
  \draw (6,3) to[R, l_=$R_3$, i>^=$I_3$] (6,0);
  \draw (6,0) -- (0,0);
\end{circuitikz}
\end{center}
\begin{teilaufgaben}
  \item Rechnen Sie die Widerstände in Leitwerte um und bestimmen Sie $I_1$, $I_2$, $I_3$ mit der Stromteilerregel für Leitwerte.
  \item Berechnen Sie $I_1$ noch einmal, diesmal mit der Stromteilerregel für \emph{zwei} Widerstände. Fassen Sie dafür $R_2$ und $R_3$ zusammen.
  \item Welche Spannung liegt an den Widerständen?
  \item Welchen Widerstand $R_4$ muss man zusätzlich parallel schalten, damit durch $R_1$ nur noch $I_1 = \SI{20}{mA}$ fließen?
\end{teilaufgaben}
\end{aufgabe}

\begin{loesung}
\schritt{a) Leitwerte}
$G_1 = \frac{1}{\SI{20}{\ohm}} = \SI{50}{mS}$, $G_2 = \SI{33.3}{mS}$, $G_3 = \SI{16.7}{mS}$, \quad $G_\text{ges} = G_1+G_2+G_3 = \SI{100}{mS}$.
\leitfrage{Mit Leitwerten sieht der Stromteiler genauso aus wie der Spannungsteiler mit Widerständen. Warum?}
Weil bei Parallelschaltung dieselbe Spannung anliegt und $I_n = G_n\cdot U$ gilt -- die Ströme verhalten sich wie die Leitwerte:
\[
  I_1 = I\frac{G_1}{G_\text{ges}} = \SI{60}{mA}\cdot\frac{50}{100} = \ergebnis{\SI{30}{mA}},\quad
  I_2 = \SI{60}{mA}\cdot\frac{33{,}3}{100} = \ergebnis{\SI{20}{mA}},\quad
  I_3 = \ergebnis{\SI{10}{mA}}
\]
Kontrolle mit dem Knotensatz: $30 + 20 + 10 = 60$ \checkmark. Der kleinste Widerstand führt den größten Strom.
\schritt{b) Zwei-Widerstands-Regel}
$R_{23} = R_2\parallel R_3 = \frac{30\cdot 60}{90}\,\si{\ohm} = \SI{20}{\ohm}$. Beim Stromteiler mit Widerständen steht im Zähler der \emph{andere} Widerstand:
\[ I_1 = I\cdot\frac{R_{23}}{R_1 + R_{23}} = \SI{60}{mA}\cdot\frac{20}{40} = \ergebnis{\SI{30}{mA}} \;\checkmark \]
\schritt{c) Spannung}
$U = \frac{I}{G_\text{ges}} = \frac{\SI{60}{mA}}{\SI{100}{mS}} = \ergebnis{\SI{0.6}{V}}$ \quad (Probe: $I_1R_1 = \SI{30}{mA}\cdot\SI{20}{\ohm} = \SI{0.6}{V}$ \checkmark)
\schritt{d) Zusätzlicher Parallelwiderstand}
\leitfrage{Welcher Anteil von $I$ soll durch $G_1$ fließen, und welchen Gesamtleitwert erfordert das?}
$I_1/I = 20/60 = 1/3 = G_1/G_\text{ges}' \;\Rightarrow\; G_\text{ges}' = 3G_1 = \SI{150}{mS}$.
\[ G_4 = G_\text{ges}' - G_\text{ges} = \SI{150}{mS} - \SI{100}{mS} = \SI{50}{mS} \quad\Rightarrow\quad R_4 = \ergebnis{\SI{20}{\ohm}} \]
\typfehler{Stromteiler mit Widerständen wie einen Spannungsteiler aufschreiben ($I_1 = I\cdot R_1/(R_1+R_{23})$) -- dann bekäme der \emph{größere} Widerstand den größeren Strom, was offensichtlich falsch ist. Mit Leitwerten passiert dieser Fehler nicht.}
\end{loesung}

\end{document}
